% Insertion fragment: place after literal_domain_quotient_addendum.tex.
% Existing algebraic theorem labels verified directly on 5 September 2026:
% sgd:parametrization and sgd:literal-addon:quotient.
% New analytic proofs below are complete and use no symbolic/numerical run.
% Audited source SHA-256 values:
% scalar_gauss_domain.tex:
% 7974E5D3501D36F232AE0CB6D68A1161DB97102D78945B40ABA2D323FD27CDCA
% literal_domain_quotient_addendum.tex:
% B0B7AA97671BACBF66B95D3AB71E80BE69CAA203C4EAA554BF32267693DD19A2
% REVIEW.md:
% 71778E75F8E1BE1C9B5AFB0B9D38701E53F3014402227495121D3601A98BF976
% No preamble or document environment. All new labels: sgd:smooth-addon:.

\section{Smooth literal inverse charts and all positive-coupling fibres}
\label{sgd:smooth-addon:section}

The exact algebraic bijections of
Theorems~\ref{sgd:parametrization} and
\ref{sgd:literal-addon:quotient} admit a smooth refinement.
We prove the parameter domain's convexity, give the phase quotient
its explicit smooth coordinates, and construct smooth inverse maps
from the literal matrix entries. We then give a global product
description in which the first continuous coordinate is the unchanged
original coupling $t=g^2$. The topology of every fixed-coupling fibre
will follow from explicit inverse formulas.

\subsection{Retained objects and the exact parameter domain}

Keep the literal matrices and constants
\[
 H_0=\begin{pmatrix}0&0&0\\0&-1/2&1/(2\sqrt2)\\
 0&1/(2\sqrt2)&1/2\end{pmatrix},\quad
 q=\frac1{2\sqrt2}\begin{pmatrix}1\\0\\0\end{pmatrix},\quad
 A_0=\sqrt2\begin{pmatrix}0&-1&0\\1&0&0\\0&0&0\end{pmatrix},
\]
\[
 m^2=\Gamma=1,\qquad \kappa=6,\qquad
 F_0=[A_0,[A_0,H_0]]=-2H_0+\operatorname{diag}(-2,1,1),
\]
and set
\[
 r=\frac{\sqrt6}{4},\quad\delta=\frac1{2r}=\sqrt{\frac23},
 \quad x=\frac12,\quad d_0=\frac1{2\sqrt2},\quad
 U=\operatorname{diag}\left(
 1,\frac1{\sqrt{2r(r+x)}}
 \begin{pmatrix}d_0&-(r+x)\\r+x&d_0\end{pmatrix}\right).
\]
Thus $U^*U=I_3$, $\det U=1$, $Uq=q$, and
$U^*H_0U=\operatorname{diag}(0,r,-r)$, as proved in
Lemma~\ref{sgd:literal-addon:temporal}.
Write
\[
 P=\operatorname{diag}(1,0,0),\qquad
 X(v)=\begin{pmatrix}0&-v^*\\v&0_2\end{pmatrix},\qquad
 \mathfrak p=X(\mathbb C^2)\subset\mathfrak{su}(3).
\]
The original real ambient manifold for this selected spatial class is
\[
 \mathcal E=(0,\infty)\times(\mathbb R\setminus\{0\})
             \times\mathfrak{su}(3)\times\mathfrak p\times\mathfrak p.
\]
It has real dimension $1+1+8+4+4=18$.
For $Y_{\rm tf}=Y-(\operatorname{tr}Y)I_3/3$, put
\[
 \mathcal I_{H_0}Y=([H_0,[H_0,Y]]+qq^*Y+Yqq^*)_{\rm tf},
 \qquad \mathcal K_{Z_1,Z_2}Y=-\sum_{j=1}^2[Z_j,[Z_j,Y]].
\]
The subset $\mathcal S_{\rm sg}\subset\mathcal E$ is exactly the
simultaneous zero set
\begin{align*}
 [C_0,[C_0,H_0]]
     -\sum_j[Z_{j,0},[Z_{j,0},H_0]]-F_0&=0,\\
 (C_0^2-Z_{1,0}^2-Z_{2,0}^2)q+2q&=0,\\
 (t\mathcal I_{H_0}+2\mathcal K_{Z_{1,0},Z_{2,0}})C_0
                   -4\omega[Z_{1,0},Z_{2,0}]&=0.
\end{align*}
Neither additional spatial residual is included in this definition.
The topology of $\mathcal S_{\rm sg}$ initially means its subspace
topology in $\mathcal E$; its embedded smooth structure will be proved.

For real coordinates $(u,v_0,y)$ define
\[
 h=\delta^2+y^2,\qquad m=\delta+iy,\qquad
 M=\begin{pmatrix}u&m\\\bar m&v_0\end{pmatrix},\qquad
 d=uv_0-h,\qquad c=4v_0-5u,
\]
\[
 \mathcal D=\{(u,v_0,y)\in\mathbb R^3:u>0,\ d>0,\ c>0\}.
\]
Here $h$ is a scalar, distinct from the two factor entries
$h_\pm$ used below.

\begin{lemma}[Convexity of the entire continuous parameter domain]
\label{sgd:smooth-addon:convex}
The set $\mathcal D$ is a nonempty open convex subset of
$\mathbb R^3$. In particular it is connected and has real dimension three.
\end{lemma}
\begin{proof}
The inequalities are strict inequalities for continuous functions,
so $\mathcal D$ is open. For $u>0$ and $(a,b)\in\mathbb C^2$,
the exact quadratic-form identity is
\[
 \begin{pmatrix}\bar a&\bar b\end{pmatrix}
 M\begin{pmatrix}a\\b\end{pmatrix}
 =u\left|a+\frac m u b\right|^2+\frac d u|b|^2.
\]
It proves $M>0$ when $u>0,d>0$. Conversely, $M>0$ gives
$u>0$ by the first coordinate vector, and evaluation at
$(a,b)=(-m/u,1)$ gives $d/u>0$.
Thus the conditions $u>0,d>0$ are exactly positive definiteness
of the Hermitian matrix in the affine slice $\operatorname{Re}M_{12}=\delta$.

If $M_0,M_1$ are positive definite and $0\le\lambda\le1$,
then for every nonzero $z\in\mathbb C^2$,
\[
 z^*((1-\lambda)M_0+\lambda M_1)z
 =(1-\lambda)z^*M_0z+\lambda z^*M_1z>0.
\]
The affine slice is preserved by this combination, and the additional
linear inequality $4v_0-5u>0$ is preserved as well.
The coordinate map $(u,v_0,y)\mapsto M$ is affine and injective,
so these facts prove convexity in the original real coordinates.
Finally $(u,v_0,y)=(\delta,2\delta,0)$ has
$u=\delta>0$, $d=\delta^2>0$, and $c=3\delta>0$.
It supplies a point and, by openness, a three-dimensional neighbourhood.
Straight segments inside a convex set join any two points,
proving connectedness.
\end{proof}

\subsection{The phase quotient with every original phase retained}

Write $\mathbb T=U(1)=\{z\in\mathbb C:|z|=1\}$.
On the original three phases, in the order $(\zeta,p_+,p_-)$,
define
\begin{equation}
 \Psi:\mathbb T^3\longrightarrow\mathbb T^3,\qquad
 \Psi(\zeta,p_+,p_-)
       =\left(\zeta^2,\frac{p_+}{\zeta},\frac{p_-}{\zeta}\right)
       =(\eta,\alpha,\beta).
 \label{sgd:smooth-addon:phase-map}
\end{equation}

\begin{lemma}[Exact smooth phase quotient]
\label{sgd:smooth-addon:phase-quotient}
The map $\Psi$ is a surjective smooth group homomorphism with kernel
$\{(1,1,1),(-1,-1,-1)\}$.
It induces a diffeomorphism
\[
 \mathbb T^3/\{(1,1,1),(-1,-1,-1)\}\ \longrightarrow\ \mathbb T^3.
\]
The quotient smooth charts have explicit local inverses retaining
both original square-root representatives.
\end{lemma}
\begin{proof}
Products, inverses on $\mathbb T$, and the square map are smooth;
the displayed formulas commute with coordinatewise multiplication.
An element is in the kernel precisely when $\zeta^2=1$ and
$p_+=p_-=\zeta$, giving exactly the two listed elements.
Given any $(\eta,\alpha,\beta)\in\mathbb T^3$, choose either root
$\zeta^2=\eta$ and set $p_+=\zeta\alpha$, $p_-=\zeta\beta$.
This gives all preimages, so $\Psi$ is onto and each fibre is
exactly one simultaneous-sign pair.

Here are the smooth local charts, rather than only a set quotient.
Fix $\eta_0\in\mathbb T$ and one root $\zeta_0^2=\eta_0$.
On the open arc $\eta=\eta_0e^{i\theta}$ with $-\pi<\theta<\pi$,
the two lifts are exactly
\[
 \zeta_\pm=\pm\zeta_0e^{i\theta/2},\qquad
 p_{+,\pm}=\zeta_\pm\alpha,\qquad
 p_{-,\pm}=\zeta_\pm\beta.
\]
Each formula is smooth. The two root arcs are disjoint and open
in the preimage, and each maps diffeomorphically to the displayed
arc times $\mathbb T^2$, using \eqref{sgd:smooth-addon:phase-map}
as its inverse. Their simultaneous-sign interchange identifies
exactly these two charts.
The quotient topology agrees with these charts: the quotient
projection is open because the inverse image of the image of any
open set is that set united with its simultaneous-sign translate.
On overlapping arcs the quotient coordinates are the same
$(\eta,\alpha,\beta)$, so their transitions are the ordinary smooth
torus transitions. These charts prove the asserted diffeomorphism.
They do not select a global square root or discard either representative.
\end{proof}

\subsection{Global formulas for literal matrices and their smooth inverses}

We record all the algebraic parameters used by the smooth map.
For $(u,v_0,y)\in\mathcal D$, set
\[
 Q=\operatorname{diag}(1,2),\qquad E=\operatorname{diag}(2/3,-5/6),
 \qquad t=\frac{72d}{c},\qquad
 T=\frac t4E-\frac12(QM+MQ).
\]
For independent signs $\epsilon,\sigma\in\{1,-1\}$ put
\[
 e=\epsilon\sqrt{2+u+v_0},\qquad
 k=\sigma\sqrt{\frac d{-\det T}},\qquad
 \omega=e/k,\qquad N=kT,\qquad A_\pm=M\pm N.
\]
Theorem~\ref{sgd:parametrization} proves on this entire domain that
$-\det T>0$, that $A_\pm$ are positive semidefinite of rank one,
and that $h_\pm=(A_\pm)_{11}>0$. Consequently
\[
 a_\pm=\begin{pmatrix}\sqrt{h_\pm}\\
                 (A_\pm)_{21}/\sqrt{h_\pm}\end{pmatrix}
\]
are defined everywhere in $\mathcal D$ for both signs.

For $(\eta,\alpha,\beta)\in\mathbb T^3$ put
\[
 v^\flat=\frac{\alpha a_++\beta a_-}{2},\qquad
 w^\flat=\frac{i(\alpha a_+-\beta a_-)}{2},\qquad
 L(\eta)=\frac1{\sqrt2}\begin{pmatrix}-1&-1\\\eta&-\eta\end{pmatrix},
 \qquad B^\flat(\eta)=U\operatorname{diag}(1,L(\eta)).
\]
The superscript $\flat$ names these exact intermediate coordinates;
it does not denote a metric dual or a change of magnitude.
Multiplication gives
\[
 L(\eta)^*L(\eta)=I_2,\qquad
 \det L(\eta)=\eta,\qquad \det B^\flat(\eta)=\eta.
\]
In particular these auxiliary matrices are unitary with their
displayed determinant; they are not asserted to have determinant one.
The precise relation to the original determinant-one frame is,
for either root $\zeta^2=\eta$,
\[
 \gamma=-\zeta^{-2},\qquad
 W(\zeta)=\frac1{\sqrt2}
 \begin{pmatrix}-\zeta^{-1}&-\zeta^{-1}\\\zeta&-\zeta\end{pmatrix},
 \qquad L(\eta)=\zeta W(\zeta).
\]
If $p_+=\zeta\alpha$, $p_-=\zeta\beta$, the original vectors satisfy
$v=\zeta v^\flat$ and $w=\zeta w^\flat$.
Thus every phase and every determinant factor is retained.

Let
\[
 \mathcal X=\{1,-1\}^2\times\mathcal D\times\mathbb T^3
\]
have coordinates $(\epsilon,\sigma,u,v_0,y,\eta,\alpha,\beta)$.
Define $\mathcal F^\flat:\mathcal X\to\mathcal E$ by
\begin{equation}
 \begin{gathered}
 \mathcal F^\flat(\epsilon,\sigma,u,v_0,y,\eta,\alpha,\beta)
       =(t,\omega,C_0,Z_{1,0},Z_{2,0}),\\
 C_0=B^\flat(\eta)\bigl(ie\operatorname{diag}(1,0,-1)\bigr)
                                    B^\flat(\eta)^*,\\
 \begin{aligned}
 Z_{1,0}&=B^\flat(\eta)X(v^\flat)B^\flat(\eta)^*,\\
 Z_{2,0}&=B^\flat(\eta)X(w^\flat)B^\flat(\eta)^*.
 \end{aligned}
 \end{gathered}
 \label{sgd:smooth-addon:literal-forward}
\end{equation}

\begin{theorem}[Four smooth connected components in literal coordinates]
\label{sgd:smooth-addon:literal-diffeomorphism}
The image of $\mathcal F^\flat$ is exactly $\mathcal S_{\rm sg}$,
and $\mathcal F^\flat$ is a diffeomorphism onto an embedded
six-real-dimensional submanifold of $\mathcal E$.
The submanifold has exactly four connected components, indexed by
\[
 \epsilon=\operatorname{sign}e,\qquad
 \sigma=\operatorname{sign}(e/\omega).
\]
Each component is diffeomorphic to $\mathcal D\times\mathbb T^3$.
The smooth structure and inverse charts refer to the literal
matrix subset; no quotient of literal configurations by a physical
gauge action is being taken.
\end{theorem}
\begin{proof}
First verify exactly which literal matrices the displayed formulas give.
Let $D_\zeta=\operatorname{diag}(1,\zeta I_2)$ and let
$B(\zeta)=U\operatorname{diag}(1,W(\zeta))$ be the original frame.
Then
\begin{align*}
 B^\flat(\eta)&=B(\zeta)D_\zeta,\\
 D_\zeta\bigl(ie\operatorname{diag}(1,0,-1)\bigr)D_\zeta^*
       &=ie\operatorname{diag}(1,0,-1).
\end{align*}
\[
 D_\zeta X(v^\flat)D_\zeta^*=X(\zeta v^\flat)=X(v),
 \qquad D_\zeta X(w^\flat)D_\zeta^*=X(w).
\]
Thus \eqref{sgd:smooth-addon:literal-forward} gives exactly the same
literal tuple as the original map for each of its two phase lifts.
Lemma~\ref{sgd:smooth-addon:phase-quotient} and the already-proved
literal quotient bijection, Theorem~\ref{sgd:literal-addon:quotient},
therefore show that its image is $\mathcal S_{\rm sg}$ and that
$\mathcal F^\flat$ is bijective.

Here is its inverse directly in literal matrix entries.
Given a tuple in $\mathcal S_{\rm sg}$, compute
\begin{equation}
 C_d=U^*C_0U,\qquad
 e=-i(C_d)_{11},\qquad z=(C_d)_{23},\qquad
 \eta=-\frac{2i\bar z}{e}.
 \label{sgd:smooth-addon:literal-temporal-inverse}
\end{equation}
The temporal classification proves that $e$ is real and nonzero,
$|z|=|e|/2$, and hence $|\eta|=1$.
The sign is checked by multiplying the lower temporal block in
\eqref{sgd:smooth-addon:literal-forward}:
\[
 (C_d)_{\rm lower}
 =-\frac{ie}{2}\begin{pmatrix}1&\bar\eta\\\eta&1\end{pmatrix}.
\]
In particular $z=-ie\bar\eta/2$, so
$-2i\bar z/e=\eta$ exactly.

Define the two matrices and their entire lower first columns by
\[
 Z_j^\flat=B^\flat(\eta)^*Z_{j,0}B^\flat(\eta),\qquad
 v^\flat=\begin{pmatrix}(Z_1^\flat)_{21}\\(Z_1^\flat)_{31}\end{pmatrix},
 \qquad
 w^\flat=\begin{pmatrix}(Z_2^\flat)_{21}\\(Z_2^\flat)_{31}\end{pmatrix}.
\]
Then recover
\[
 M=v^\flat(v^\flat)^*+w^\flat(w^\flat)^*,\qquad
 N=i\bigl(v^\flat(w^\flat)^*-w^\flat(v^\flat)^*\bigr),
\]
\begin{equation}
 \begin{gathered}
 u=M_{11},\quad v_0=M_{22},\quad y=\operatorname{Im}M_{12},
 \qquad \epsilon=\operatorname{sign}e,\quad
 \sigma=\operatorname{sign}(e/\omega),\\
 h_\pm=(M\pm N)_{11},\qquad
 \alpha=\frac{(v^\flat-iw^\flat)_1}{\sqrt{h_+}},\qquad
 \beta=\frac{(v^\flat+iw^\flat)_1}{\sqrt{h_-}}.
 \end{gathered}
 \label{sgd:smooth-addon:literal-spatial-inverse}
\end{equation}
Indeed multiplication by the common unit $\zeta$ preserves both
$M$ and $N$. The original phase inverse therefore gives
$\alpha=p_+/\zeta$, $\beta=p_-/\zeta$ in these formulas, with
$h_\pm>0$, $|\alpha|=|\beta|=1$.
The algebraic inverse proves that all recovered coordinates lie in
their exact domains. The displayed formulas recover the original
parameter from \eqref{sgd:smooth-addon:literal-forward}, and the
forward map recovers the original literal tuple from them.

All forward functions are real smooth on each sign sector:
$c>0$, $d>0$, $-\det T>0$, $2+u+v_0>0$, and $h_\pm>0$
are precisely the denominators and square-root arguments appearing
above, while $k\ne0$. Matrix products and conjugates are real smooth.
In the inverse, $e\ne0$ and $h_\pm>0$ ensure smoothness of every
non-discrete formula. The signs are locally constant.
The inverse formulas in particular give a continuous inverse for
the subspace topology on $\mathcal S_{\rm sg}$.

We check embeddedness, rather than just transferring an abstract
smooth structure through a bijection.
At any one image point choose open torus arcs around
$\eta,\alpha,\beta$, with their real angle coordinates.
The inverse formulas extend as real smooth functions to an open
neighbourhood in $\mathcal E$ as follows. Keep $e\ne0$ and the
computed $h_\pm>0$ by taking a sufficiently small neighbourhood.
The same matrix formulas define the complex quantities
$\eta,\alpha,\beta$ there, even when the tuple is off the zero set.
For a unit base point $\eta_0$, on the complex neighbourhood
$\operatorname{Re}(\eta/\eta_0)>0$, its local angle is the smooth
real function
\[
 \arctan\!\left(
 \frac{\operatorname{Im}(\eta/\eta_0)}
      {\operatorname{Re}(\eta/\eta_0)}\right).
\]
Use the corresponding expressions about $\alpha_0,\beta_0$.
On the zero set these are exactly its three local torus coordinates.
Together with $u,v_0,y$, they give a real smooth six-coordinate
map from an ambient neighbourhood whose composition with
$\mathcal F^\flat$ on the parameter chart is the identity.
The derivative of that composition is the identity, so the
derivative of $\mathcal F^\flat$ has rank six at every point.

For clarity, the standard inverse function theorem now applies with
its full-rank hypothesis verified: choose six ambient linear
coordinate projections whose composite with this rank-six
parametrization has invertible derivative. Locally the image is a
smooth graph over those six coordinates.
The continuous inverse already proved prevents other parameter
patches from contributing additional image points arbitrarily near
this point outside its chart: the inverse image of the chart is
open in the subspace topology. Thus the graph is a neighbourhood
of the point in the whole subset $\mathcal S_{\rm sg}$.
This proves it is an embedded six-dimensional submanifold and that
the displayed inverse is smooth for its embedded structure.

It remains to count components exactly. By
Lemma~\ref{sgd:smooth-addon:convex}, $\mathcal D$ is path connected.
Each circle is path connected: for chosen real arguments of two
unit endpoints, interpolating those arguments gives an explicit
path on the circle. Products of these paths prove $\mathbb T^3$
is path connected. Hence each of the four sign sectors in
$\mathcal X$ is nonempty and connected.
On the literal domain the two sign functions in
\eqref{sgd:smooth-addon:literal-spatial-inverse} are continuous
with values in $\{1,-1\}$ because $e,\omega$ never vanish.
They cannot change on a connected subset, so distinct sectors
cannot belong to the same component. This proves there are
exactly four connected components.

Finally the only quotient used to construct the coordinates was
the two-element kernel of the parameter map.
Its two representatives give identical literal matrices, as the
first displayed equalities proved. Distinct literal tuples remain
distinct under the bijection. No further gauge equivalence or
spatial-pair equivalence has been imposed.
\end{proof}

\subsection{The original coupling is a surjective submersion}

On $\mathcal D$ the original coupling function is
\[
 t(u,v_0,y)=\frac{72(uv_0-\delta^2-y^2)}{4v_0-5u}
          =\frac{72d}{c}.
\]

\begin{theorem}[Exact differential and global coupling coordinates]
\label{sgd:smooth-addon:coupling}
The function $t:\mathcal D\to(0,\infty)$ is a smooth surjective
submersion. Its exact partial derivatives are
\begin{equation}
 \begin{gathered}
 \frac{\partial t}{\partial u}
 =\frac{72(4v_0^2-5h)}{c^2},\qquad
 \frac{\partial t}{\partial v_0}
 =\frac{72(4h-5u^2)}{c^2},\\
 \frac{\partial t}{\partial y}=-\frac{144y}{c},
 \qquad h=\delta^2+y^2.
 \end{gathered}
 \label{sgd:smooth-addon:derivatives}
\end{equation}
In fact $\partial t/\partial u>0$ throughout $\mathcal D$.

Define the fixed open set
\[
 \mathcal B=\{(v_0,y)\in\mathbb R^2:
        v_0>\tfrac{\sqrt5}{2}\sqrt{\delta^2+y^2}\}.
\]
The map
\begin{equation}
 \Theta:\mathcal D\longrightarrow(0,\infty)\times\mathcal B,
 \qquad \Theta(u,v_0,y)=(t(u,v_0,y),v_0,y)
 \label{sgd:smooth-addon:trivialization}
\end{equation}
is a diffeomorphism, with exact inverse
\begin{equation}
 u=\frac{4\tau v_0+72(\delta^2+y^2)}{72v_0+5\tau},
 \qquad (\tau,v_0,y)\in(0,\infty)\times\mathcal B.
 \label{sgd:smooth-addon:coupling-inverse}
\end{equation}
Moreover $\mathcal B$ is diffeomorphic to $\mathbb R^2$ by the
explicit inverse pair
\begin{equation}
 \begin{aligned}
 (y,s)&\longmapsto
 \left(\tfrac{\sqrt5}{2}\sqrt{\delta^2+y^2}+e^s,\ y\right),\\
 (v_0,y)&\longmapsto
 \left(y,\log\!\left(v_0-\tfrac{\sqrt5}{2}\sqrt{\delta^2+y^2}\right)\right).
 \end{aligned}
 \label{sgd:smooth-addon:B-chart}
\end{equation}
Here $\tau=t=g^2$ is the unchanged positive original coupling.
\end{theorem}

\begin{proof}
Since $c>0$, the quotient defining $t$ is smooth.
The derivatives $d_u=v_0$, $d_{v_0}=u$, $d_y=-2y$,
$c_u=-5$, $c_{v_0}=4$, $c_y=0$ give
\[
 \frac{\partial t}{\partial u}
 =\frac{72(v_0c+5d)}{c^2}
 =\frac{72(4v_0^2-5h)}{c^2},
 \qquad
 \frac{\partial t}{\partial v_0}
 =\frac{72(uc-4d)}{c^2}
 =\frac{72(4h-5u^2)}{c^2},
\]
and the last derivative in \eqref{sgd:smooth-addon:derivatives}
is $72(-2y)/c$.
On $\mathcal D$, $v_0=(d+h)/u>0$, $c>0$, and $d>0$.
Hence $v_0c+5d>0$, proving $\partial t/\partial u>0$.
The differential therefore has rank one everywhere, proving the
submersion assertion independently of any fibre-topology argument.
In particular its three partial derivatives cannot simultaneously
vanish. The proposed algebraic simultaneous-zero test also agrees:
the first two zero equations would give
$4v_0^2=5h$ and $4h=5u^2$, so positivity of $u,v_0$ would imply
$v_0=5u/4$, and then both $c=0$ and $d=0$, excluded from $\mathcal D$.

We now prove the exact global inverse and surjectivity.
For a point of $\mathcal D$, the identity
\[
 4v_0^2-5h=v_0c+5d>0
\]
and $v_0>0$ imply $(v_0,y)\in\mathcal B$.
Writing $\tau=72d/c>0$, multiply by the retained $c$ to obtain
\[
 \tau(4v_0-5u)=72(uv_0-h),
 \qquad
 (72v_0+5\tau)u=4\tau v_0+72h.
\]
Its denominator $72v_0+5\tau$ is strictly positive, giving
\eqref{sgd:smooth-addon:coupling-inverse}.

Conversely fix any $\tau>0$ and $(v_0,y)\in\mathcal B$, and use
\eqref{sgd:smooth-addon:coupling-inverse}. Here $h>0$, $v_0>0$,
and the numerator and denominator of $u$ are positive.
Exact multiplication, with no sign or factor removed, gives
\begin{align}
 c=4v_0-5u
 &=\frac{72(4v_0^2-5h)}{72v_0+5\tau}>0,\notag\\
 d=uv_0-h
 &=\frac{\tau(4v_0^2-5h)}{72v_0+5\tau}>0.
 \label{sgd:smooth-addon:inverse-domain}
\end{align}
Thus the recovered point is in $\mathcal D$, and
$72d/c=\tau$ exactly.
These calculations prove both compositions of
\eqref{sgd:smooth-addon:trivialization} and
\eqref{sgd:smooth-addon:coupling-inverse} are identities.
Both maps are smooth on their stated open domains.

The set $\mathcal B$ is nonempty; for example
$(v_0,y)=(\sqrt5\delta/2+1,0)$ belongs to it.
Thus every $\tau>0$ has a preimage, proving surjectivity.
Finally $\delta>0$ makes $\sqrt{\delta^2+y^2}$ smooth for all real
$y$, and the defining inequality of $\mathcal B$ makes the
logarithm argument in \eqref{sgd:smooth-addon:B-chart} strictly
positive. Exponential and logarithm are inverse smooth functions
on their stated real domains. Substitution into
\eqref{sgd:smooth-addon:B-chart} proves both compositions are
identities, establishing the diffeomorphism with $\mathbb R^2$.
\end{proof}

\subsection{Exact scope and topology of the positive-coupling level sets}

\begin{corollary}[All positive-coupling levels and their literal inverse charts]
\label{sgd:smooth-addon:levels}
In the literal embedded manifold $\mathcal S_{\rm sg}$, the
original coupling projection
\[
 \pi:\mathcal S_{\rm sg}\longrightarrow(0,\infty),\qquad
 (t,\omega,C_0,Z_{1,0},Z_{2,0})\longmapsto t
\]
is a smooth surjective submersion. There is an explicit diffeomorphism
\begin{equation}
 \mathcal S_{\rm sg}
 \ \cong\
 \{1,-1\}^2\times(0,\infty)\times\mathbb R^2\times\mathbb T^3
 \label{sgd:smooth-addon:global-product}
\end{equation}
under which $\pi$ is projection onto the displayed unchanged
$(0,\infty)$ coordinate.
For every specified original $\tau>0$, the complete fixed-coupling
level $\pi^{-1}(\tau)$ is an embedded five-real-dimensional
submanifold, with exactly four connected components.
Each component is diffeomorphic to $\mathbb R^2\times\mathbb T^3$.
These statements concern precisely the scalar--Gauss domain with
both original signs and all three phase coordinates retained.
\end{corollary}
\begin{proof}
Compose the literal diffeomorphism of
Theorem~\ref{sgd:smooth-addon:literal-diffeomorphism} with
the diffeomorphisms of
Theorem~\ref{sgd:smooth-addon:coupling}. More explicitly, given
\[
 (\epsilon,\sigma,\tau,y,s,\eta,\alpha,\beta)
 \in\{1,-1\}^2\times(0,\infty)\times\mathbb R^2\times\mathbb T^3,
\]
first put
\[
 v_0=\frac{\sqrt5}{2}\sqrt{\delta^2+y^2}+e^s,\qquad
 u=\frac{4\tau v_0+72(\delta^2+y^2)}{72v_0+5\tau}.
\]
Then apply \eqref{sgd:smooth-addon:literal-forward} with these
$u,v_0,y$ and the retained signs and phases.
Equation~\eqref{sgd:smooth-addon:inverse-domain} proves that its
coupling is exactly $\tau$ and every denominator is positive.
The inverse begins with the literal matrix formulas
\eqref{sgd:smooth-addon:literal-temporal-inverse} and
\eqref{sgd:smooth-addon:literal-spatial-inverse}, keeps
$\tau=t$ from the original tuple, and recovers
\[
 s=\log\!\left(v_0-\frac{\sqrt5}{2}\sqrt{\delta^2+y^2}\right).
\]
The preceding theorems prove that these maps are smooth, mutually
inverse, and defined on every point of their stated domains.
This proves \eqref{sgd:smooth-addon:global-product}.

In these exact product coordinates $\pi$ is the projection onto
$\tau$. Its derivative is onto the one-dimensional tangent space
of $(0,\infty)$, and every $\tau>0$ occurs.
The fixed-coordinate slice is a smooth embedded codimension-one
submanifold. Transport through the diffeomorphism makes
$\pi^{-1}(\tau)$ an embedded five-dimensional submanifold of
$\mathcal S_{\rm sg}$, and hence of the original ambient manifold
$\mathcal E$. Its inclusion into the fixed-coupling ambient
hyperplane has the same smooth structure, since that coordinate is
identically $\tau$.

The formula identifies the entire fibre, without a missing
parameter portion, with
$\{1,-1\}^2\times\mathbb R^2\times\mathbb T^3$.
The real plane is path connected by line segments; the three
circles are path connected by the explicit angle paths used
earlier. Their product is therefore connected.
The two continuous nonvanishing sign invariants prevent any
connection between distinct sign sectors. This proves both
the exact count of four components and their stated diffeomorphism
type. The fibre topology follows from these explicit inverse maps,
not merely from the fact that $\pi$ is a submersion.
\end{proof}

These manifolds are subsets of the space of ordered literal tuples.
The only phase identification was equality of two parameter
representatives producing the same tuple. No spatial equation has
been imposed, and no physical gauge quotient, frequency-sign
identification, or exchange of $Z_{1,0}$ and $Z_{2,0}$ has been made.
The construction and all level-set statements use the open domain
$t>0$; they assign no boundary fibre at $t=0$.
